\itemhead{复习题三 · 第18题（教材p.92）}{XJ-B1-C03-RV-P092-Y0375}
\label{XJ-B1-C03-RV-P092-Y0375}
{\footnotesize 湘教版教材p.92 / PDF第99页 / 复习题三 / 第18题}
\field{原题}{已知函数 $y=f(x)$ 与 $y=g(x)$ 的图象如图所示，则函数 $y=f(x)\cdot g(x)$ 的图象可能是（\quad）
\bankfig{C03-RV092}{0.88}}
\field{参考答案}{A。$f$ 偶、$g$ 奇，乘积为奇函数；在0右侧邻近处 $f>0,g<0$，乘积为负，排除其他图象。}
{\small \field{关联知识}{F04 函数三种表示法、图象与读数（教材帮p.85、86、93、94 / PDF94、95、102、103）；F08 奇偶性、对称与函数性质应用（教材帮p.113、114、115 / PDF122、123、124）}
\field{方法与易错点}{把关键点与定义域一起标到图上，逐区间读交点和上下关系。 易错：将示意图预测当事实；把两曲线交点收进严格不等式；折线漏端点。}
\field{答案依据}{答案册相关题组：PDF第24、25、26、27页（册内页23、24、25、26）；按当前题库原题逐项重算，不按旧版题号机械对套。}}
